\chapter*{ఓపెన్ లాజిక్ ప్రాజెక్ట్ గురించి}
\addcontentsline{toc}{chapter}{ఓపెన్ లాజిక్ ప్రాజెక్ట్ గురించి}


\textit{Open Logic Text} అనేది ఆచారబద్ధ అధితర్కం, ఆచారబద్ధ పద్ధతులపై
బహిరంగ మూలాలతో, సహకారంతో రూపొందుతున్న పాఠ్యపుస్తకం. ఇది మధ్యస్థ స్థాయి
నుంచి మొదలవుతుంది; అంటే ఆచారబద్ధ తర్కానికి పరిచయ పాఠ్యక్రమం పూర్తయిన
తరువాత చదివే స్థాయి. ముఖ్యంగా తత్వశాస్త్రం, కంప్యూటర్ సైన్స్
విద్యార్థుల వంటి గణితేతర పాఠకులను దృష్టిలో ఉంచుకున్నా, దాని
వివరణలు గణితపరంగా కచ్చితమైనవి.

ప్రస్తుతం చేర్చిన కొన్ని అంశాల వివరణ ఇంకా పూర్తికాకపోవచ్చు;
అనేక విభాగాలకు గణనీయమైన సవరణ అవసరం. భవిష్యత్తులో మరిన్ని
అంశాలకు పాఠ్యాన్ని విస్తరించాలని భావిస్తున్నాం. పదకోశం, అదనపు
పఠనాల జాబితా, చారిత్రక గమనికలు, చిత్రాలు, మెరుగైన వివరణలు,
ఫలితాలు తత్వశాస్త్రం, కంప్యూటర్ సైన్స్, గణితానికి ఎందుకు
సంబంధించాయో తెలిపే విభాగాలు, మరిన్ని సాధనలు, ఉదాహరణలు కూడా
చేర్చాలని భావిస్తున్నాం. మీకు దోషం కనిపించినా, సూచన ఉన్నా,
\href{https://github.com/OpenLogicProject/OpenLogic/wiki/Contributing}{ప్రాజెక్ట్ బృందానికి తెలియజేయండి}.

ఈ ప్రాజెక్ట్ బహిరంగ మూలాల స్ఫూర్తితో నడుస్తుంది. పాఠ్యం ఉచితంగా
అందుబాటులో ఉండటమే కాక, దాని LaTeX మూలాన్ని క్రియేటివ్ కామన్స్
ఆపాదన అనుమతితో ఇస్తున్నాం. తగిన మూల ఆపాదన ఇస్తే, ఎవరైనా ఈ
పనిని దింపుకోవచ్చు, ఉపయోగించవచ్చు, సవరించవచ్చు, క్రమాన్ని
మార్చవచ్చు, ఇతర రూపాల్లోకి మార్చవచ్చు, మళ్లీ పంపిణీ చేయవచ్చు.
అదనపు వివరాలకు ఓపెన్ లాజిక్ ప్రాజెక్ట్ వెబ్‌సైటు
\href{http://openlogicproject.org/}{openlogicproject.org} చూడండి.
